% =====================================================================
%  BÖLÜM 2 — Olasılık Uzayları
%  Kaynaklar: Kolmogorov 1933; Feller Vol. I Böl. I-V; Gnedenko Böl. 1
% =====================================================================

\chapter{Olasılık Uzayları}

\bolumalinti{``The theory of probability as mathematical discipline can
  and should be developed from axioms in exactly the same way as Geometry
  and Algebra.''}{Andrey Kolmogorov, \emph{Grundbegriffe der Wahrscheinlichkeitsrechnung} (1933)}

\noindent\textbf{Ön Koşullar:} Bölüm~0 (küme cebiri), Bölüm~1 (ölçüm teorisi).

\noindent\textbf{Bu Bölüm Sonraki Bölümlerde Kullanılır:}
Bölüm~3 (rasgele değişkenler),
Bölüm~4 (beklenti),
Bölüm~7 (büyük sayılar).

\noindent\textbf{Tahmini Okuma Süresi:} $\approx 8$~saat.

\vspace{0.5\baselineskip}

\begin{kazanimlar}
Bu bölümün sonunda öğrenci:
\begin{itemize}
  \item Kolmogorov aksiyomlarını ifade edebilmeli ve bu aksiyomlardan
        olasılığın temel özelliklerini (monotonluk, alt-toplamsallık,
        Poincaré formülü) türetebilmeli,
  \item Koşullu olasılığı tanımlayabilmeli, bunun bir olasılık ölçüsü
        olduğunu gösterebilmeli ve Toplam Olasılık Teoremini ispatlayabilmeli,
  \item Bağımsızlık kavramını (olaylar, olay aileleri, $\sigma$-cebirleri)
        kesin biçimde tanımlayabilmeli; ikili bağımsızlık ile karşılıklı
        bağımsızlık arasındaki farkı karşı örnekle sergileyebilmeli,
  \item Bayes teoremini ispatlayabilmeli ve önsel/sonsal olasılık yorumunu
        açıklayabilmeli,
  \item Klasik (eşit olası) olasılık uzaylarında kombinatoryal hesaplar
        yapabilmeli: permütasyon, kombinasyon, doğum günü paradoksu.
\end{itemize}
\end{kazanimlar}

\begin{motivasyon}
Bölüm~1'de ölçüm teorisinin genel çerçevesini kurduk. Bu bölümde o çerçeveyi
olasılık için özelleştiriyoruz: ölçü fonksiyonu $\mu$ yerine \emph{olasılık
ölçüsü} $\Prob$ kullanacak, $\mu(\Omega)=1$ normalizasyonu tüm yorumu değiştirecek.

Bir \emph{deneyden} söz ettiğimizde, aklımızdaki matematiksel yapı üç nesneye
sahip: mümkün tüm sonuçların kümesi $\Omega$, iyi tanımlanmış olay kavramı
$\mathcal{F}$, ve olaylara olasılık atayan $\Prob$. Kolmogorov, 1933'te bu
üçlünün tam olarak bir olasılık ölçüm uzayı $(\Omega,\mathcal{F},\Prob)$
olduğunu gösterdi.

\medskip

\textbf{Köprü kurma.} Lise matematiğinde ``eşit olası sonuçlar'' varsayımıyla
$\Prob(A) = |A|/|\Omega|$ formülünü kullandınız. Bu bölümde hem bu klasik
yaklaşımı hem de sonsuz ya da sürekli örnek uzaylı deneyleri kapsayan genel
çerçeveyi öğreneceğiz.
\end{motivasyon}

\begin{tarihselnot}
Jacob Bernoulli 1713'te (ölümünden sonra yayınlanan \emph{Ars Conjectandi}),
Pierre-Simon Laplace 1812'de (\emph{Théorie Analytique des Probabilités})
olasılığı sayısal biçimde ele aldılar; fakat her ikisi de yalnızca sonlu
eşit-olası uzaylara dayanıyordu. 19.~yüzyılda sonsuz seri ve integrallerle
olasılık hesapları yapıldıkça tutarsızlıklar ortaya çıktı (\emph{Bertrand
paradoksu}, 1888). Kolmogorov 1933'te Lebesgue ölçüm teorisini olasılığa
uyguladı ve bu paradoksları çözdü: olasılık artık bir ölçüm uzayıdır.
\end{tarihselnot}


% =====================================================================
\section{Kolmogorov Aksiyomları ve Temel Özellikler}
% =====================================================================

\begin{tanim}[Olasılık Uzayı]\label{def:olasilikuzayi}
Üç bileşenden oluşan $(\Omega, \mathcal{F}, \Prob)$ üçlüsüne
\textbf{olasılık uzayı} \emph{(probability space)} denir; burada:
\begin{itemize}
  \item $\Omega$: \textbf{örnek uzay} \emph{(sample space)} — deneyin tüm
        olası sonuçlarının kümesi; elemanlarına $\omega$ denir ve
        \textbf{sonuç} \emph{(outcome, elementary event)} adını alır.
  \item $\mathcal{F}$: $\Omega$ üzerinde bir \textbf{$\sigma$-cebiri}
        (Tanım~\ref{def:sigmacebiri}) — elemanlarına \textbf{olay} \emph{(event)} denir.
  \item $\Prob: \mathcal{F} \to [0,1]$: \textbf{olasılık ölçüsü}
        \emph{(probability measure)} — Kolmogorov'un üç aksiyomunu sağlar:
\end{itemize}
\begin{enumerate}[label=\textbf{(K\arabic*)}]
  \item $\Prob(A) \geq 0$ her $A \in \mathcal{F}$ için,
  \item $\Prob(\Omega) = 1$,
  \item $A_1, A_2, \ldots \in \mathcal{F}$ ikili ayrık ise
        $\Prob\!\left(\bigcup_{n=1}^\infty A_n\right) = \sum_{n=1}^\infty \Prob(A_n)$.
\end{enumerate}
\end{tanim}

\begin{notkutu}
Bölüm~1, Tanım~\ref{def:olasilikolucu} ile bu tanım özdeştir: olasılık ölçüsü,
$\Prob(\Omega)=1$ şartını sağlayan bir ölçüdür. Aksiyom (K3),
$\sigma$-toplamsallığın olasılık diline tercümesidir.
\end{notkutu}

\begin{ornek}[Olasılık uzayı örnekleri]\label{ex:olasilikorn}
\begin{enumerate}[label=(\alph*)]
  \item \textbf{Adil zar:} $\Omega = \{1,2,3,4,5,6\}$,
        $\mathcal{F} = \mathcal{P}(\Omega)$,
        $\Prob(\{k\}) = 1/6$ her $k$ için. Sayma ölçüsünün
        $1/|\Omega|$ ile ölçeklenmesi.

  \item \textbf{Düzgün dağılım $[0,1]$ üzerinde:} $\Omega = [0,1]$,
        $\mathcal{F} = \mathcal{B}([0,1])$,
        $\Prob = \lambda\!\restriction_{[0,1]}$ (Lebesgue ölçüsü).
        $\Prob((a,b)) = b-a$ her $0 \leq a < b \leq 1$ için.

  \item \textbf{Sonsuz yazı-tura:} $\Omega = \{0,1\}^{\mathbb{N}}$
        (sonsuz ikili diziler), $\mathcal{F}$ silindir kümelerin ürettiği
        $\sigma$-cebiri. Kolmogorov'un uzatma teoremiyle $\Prob$ inşa edilir
        (Bölüm~7'de ele alınacak).
\end{enumerate}
\end{ornek}

\begin{teorem}[Olasılığın Temel Özellikleri]\label{thm:olasiliktemel}
$(\Omega, \mathcal{F}, \Prob)$ bir olasılık uzayı, $A, B, A_n \in \mathcal{F}$.
\begin{enumerate}[label=(\roman*)]
  \item $\Prob(\emptyset) = 0$.
  \item \textbf{Sonlu toplamsallık:} $A_1,\ldots,A_n$ ikili ayrık $\Rightarrow$
        $\Prob\!\bigl(\bigcup_{k=1}^n A_k\bigr) = \sum_{k=1}^n \Prob(A_k)$.
  \item \textbf{Tümleyen:} $\Prob(A^c) = 1 - \Prob(A)$.
  \item \textbf{Monotonluk:} $A \subseteq B \Rightarrow \Prob(A) \leq \Prob(B)$.
  \item \textbf{Fark:} $A \subseteq B \Rightarrow \Prob(B \setminus A) = \Prob(B) - \Prob(A)$.
  \item \textbf{Toplam:} $\Prob(A \cup B) = \Prob(A) + \Prob(B) - \Prob(A \cap B)$.
  \item \textbf{Alt-toplamsallık (Boole):}
        $\Prob\!\bigl(\bigcup_n A_n\bigr) \leq \sum_n \Prob(A_n)$.
  \item \textbf{Aşağıdan süreklilik:} $A_n \nearrow A \Rightarrow \Prob(A_n) \nearrow \Prob(A)$.
  \item \textbf{Yukarıdan süreklilik:} $A_n \searrow A \Rightarrow \Prob(A_n) \searrow \Prob(A)$.
\end{enumerate}
\end{teorem}

\begin{proof}
Tümü Bölüm~1, Teorem~\ref{thm:olcuozellik}'den $\mu = \Prob$ ile alınır.
Ek olarak: (iii) $A \cup A^c = \Omega$ ayrık; (K2)'den $1 = \Prob(A)+\Prob(A^c)$.
(vi) $A \cup B = A \cup (B \setminus A)$ ayrık; (ii) + (v).
Yukarıdan süreklilik için $\Prob(\Omega)=1<\infty$; Teorem~\ref{thm:olcuozellik}(v)
uygulanabilir.
\end{proof}

\begin{notkutu}
(ix) için $\Prob(A_1) < \infty$ yerine $\Prob(\Omega) = 1 < \infty$ her zaman
sağlandığından, olasılık uzayında yukarıdan süreklilik \emph{koşulsuz} geçerlidir.
Bu, genel ölçüm teorisine kıyasla önemli bir kolaylıktır.
\end{notkutu}

\begin{teorem}[Poincaré Dahil-Hariç Formülü]\label{thm:poincare}
Sonlu aile $A_1, \ldots, A_n \in \mathcal{F}$ için:
\[
  \Prob\!\left(\bigcup_{k=1}^n A_k\right)
  = \sum_{k} \Prob(A_k)
  - \sum_{k < l} \Prob(A_k \cap A_l)
  + \sum_{k < l < m} \Prob(A_k \cap A_l \cap A_m)
  - \cdots
  + (-1)^{n-1} \Prob(A_1 \cap \cdots \cap A_n).
\]
\end{teorem}

\begin{proof}
\textit{Fikir:} Bir $\omega \in \bigcup_k A_k$, tam olarak $r$ adet $A_k$'da
bulunsun. Bu $\omega$'nın formülün sağ tarafına katkısını hesaplayın.

\medskip

$\omega$'nın sol tarafa katkısı $1$'dir. Sağ tarafta: $\omega$, $\binom{r}{1}$
adet tekli kümede, $\binom{r}{2}$ adet ikili kesişimde, \ldots yer alır.
Katkı:
\[
  \binom{r}{1} - \binom{r}{2} + \binom{r}{3} - \cdots + (-1)^{r-1}\binom{r}{r}
  = 1 - (1-1)^r = 1.
\]
(Binom açılımından $\sum_{k=0}^r (-1)^k \binom{r}{k} = 0$.)
$\omega \notin \bigcup_k A_k$ ise her iki taraftaki katkı $0$'dır.
\end{proof}

\begin{ornek}[Poincaré --- iki ve üç olay]\label{ex:poincare}
\textbf{İki olay:}
$\Prob(A \cup B) = \Prob(A) + \Prob(B) - \Prob(A \cap B)$ (Teorem~\ref{thm:olasiliktemel}(vi)).

\textbf{Üç olay:}
$\Prob(A \cup B \cup C) = \Prob(A) + \Prob(B) + \Prob(C)
- \Prob(A \cap B) - \Prob(A \cap C) - \Prob(B \cap C)
+ \Prob(A \cap B \cap C)$.
\end{ornek}

\begin{mikroozet}[§2.1 özeti]
Olasılık uzayı $(\Omega,\mathcal{F},\Prob)$: örnek uzay + $\sigma$-cebiri (olaylar) +
olasılık ölçüsü. Aksiyomlardan: tümleme, monotonluk, süreklilik, Poincaré formülü.
Olasılık uzayında yukarıdan süreklilik koşulsuzdur.
\end{mikroozet}


% =====================================================================
\section{Neredeyse Kesin Olaylar ve Sıfır Olasılıklı Olaylar}
% =====================================================================

\begin{tanim}[Neredeyse Kesin Olay]\label{def:neredeyse_kesin}
$A \in \mathcal{F}$:
\begin{itemize}
  \item \textbf{neredeyse kesin (n.k.)} \emph{(almost surely, a.s.)} adını alır;
        eğer $\Prob(A) = 1$.
  \item \textbf{sıfır olasılıklı} \emph{(null event)} adını alır;
        eğer $\Prob(A) = 0$.
  \item $A^c$ sıfır olasılıklı ise $A$ neredeyse kesindir ve $A^c$'ye
        \textbf{ihmal edilebilir olay} \emph{(negligible event)} denir.
\end{itemize}
\end{tanim}

\begin{onerme}[Sayılabilir birleşim n.k.]\label{prop:nk_birlesim}
$A_n$ n.k. her $n$ için ise $\bigcap_{n=1}^\infty A_n$ n.k.'dir.
\end{onerme}

\begin{proof}
$\Prob(A_n^c) = 0$ her $n$. Alt-toplamsallık:
$\Prob\!\bigl(\bigcup_n A_n^c\bigr) \leq \sum_n \Prob(A_n^c) = 0$.
De~Morgan: $\left(\bigcap_n A_n\right)^c = \bigcup_n A_n^c$; $\Prob(\text{tümleyen}) = 0$;
yani $\Prob(\bigcap_n A_n) = 1$.
\end{proof}

\begin{hataanalizi}[Sıfır olasılık $\neq$ imkânsız]
$\Omega = [0,1]$, $\Prob = \lambda$. Herhangi bir $x \in [0,1]$ için
$\Prob(\{x\}) = \lambda(\{x\}) = 0$. Ama $\{x\}$ gerçekleşebilir
(o sonuç fiilen mümkündür); yalnızca Lebesgue ölçüsü sıfır atamıştır.
``Sıfır olasılık = imkânsız'' yerine ``sıfır olasılık = ihmal edilebilir'' demek daha doğrudur.
\end{hataanalizi}

\begin{tanim}[Tam Olasılık Uzayı]\label{def:tamuzay}
$(\Omega,\mathcal{F},\Prob)$ \textbf{tam} \emph{(complete)} adını alır; eğer
$\Prob(N)=0$ ve $A \subseteq N$ ise $A \in \mathcal{F}$ (ve $\Prob(A)=0$).
Lebesgue ölçüm uzayı tam; Borel kümeler üzerinde Lebesgue ölçüsü tam değil.
\end{tanim}


% =====================================================================
\section{Koşullu Olasılık}
% =====================================================================

\begin{kesif}[Koşullu olasılık sezgisi]
Adil bir zar atılıyor. ``Çift sayı geldi'' ($B$) bilgisiyle, ``6 geldi'' ($A$)
olayının olasılığı nedir? Sezginizi yazın, ardından Tanım~\ref{def:kos_olasilik}'i okuyun.
\ogrencialani[1.5cm]
\end{kesif}

\begin{tanimgerekce}
$B$ olduğu biliniyorsa örnek uzayımız $\Omega$ yerine $B$ olur. $A$'nın bu
küçülmüş uzaydaki payı $\Prob(A \cap B) / \Prob(B)$'dir. Payda normalizasyonu
($B$'nin ölçüsünü 1 yapıyor) sağlıyor.
\end{tanimgerekce}

\begin{tanim}[Koşullu Olasılık]\label{def:kos_olasilik}
$(\Omega,\mathcal{F},\Prob)$ olasılık uzayı, $B \in \mathcal{F}$, $\Prob(B) > 0$.
$B$'ye koşullu $A$'nın olasılığı:
\[
  \Prob(A \mid B) \coloneqq \frac{\Prob(A \cap B)}{\Prob(B)}, \quad A \in \mathcal{F}.
\]
\end{tanim}

\begin{teorem}[Koşullu Olasılık Bir Olasılık Ölçüsüdür]\label{thm:kos_olcu}
$\Prob(B) > 0$ ise $A \mapsto \Prob(A \mid B)$ bir olasılık ölçüsüdür;
yani $(\Omega, \mathcal{F}, \Prob(\cdot \mid B))$ bir olasılık uzayıdır.
\end{teorem}

\begin{proof}
(K1) $\Prob(A \cap B) \geq 0$; $\Prob(B) > 0$; dolayısıyla $\Prob(A \mid B) \geq 0$.

(K2) $\Prob(\Omega \mid B) = \Prob(\Omega \cap B)/\Prob(B) = \Prob(B)/\Prob(B) = 1$.

(K3) $(A_n)$ ikili ayrık; $(A_n \cap B)$ da ikili ayrık:
\[
  \Prob\!\left(\bigcup_n A_n \,\Big|\, B\right)
  = \frac{\Prob\!\left(\bigcup_n (A_n \cap B)\right)}{\Prob(B)}
  = \frac{\sum_n \Prob(A_n \cap B)}{\Prob(B)}
  = \sum_n \Prob(A_n \mid B).  \qed
\]
\end{proof}

\begin{teorem}[Toplam Olasılık Teoremi]\label{thm:toplam_olasilik}
$(B_n)_{n=1}^N$ ($N \leq \infty$) ikili ayrık, $\bigcup_n B_n = \Omega$,
$\Prob(B_n) > 0$ olsun. Her $A \in \mathcal{F}$ için:
\[
  \Prob(A) = \sum_{n=1}^N \Prob(A \mid B_n)\, \Prob(B_n).
\]
\end{teorem}

\begin{proof}
$A = \bigcup_n (A \cap B_n)$ ayrık birleşim; $\sigma$-toplamsallık:
$\Prob(A) = \sum_n \Prob(A \cap B_n) = \sum_n \Prob(A \mid B_n)\Prob(B_n)$.
\end{proof}

\begin{ornek}[Toplam Olasılık --- fabrika sorunu]\label{ex:fabrika}
Üç fabrika sırasıyla üretimin $\%50$, $\%30$, $\%20$'sini yapıyor ve
hatalı oranları sırasıyla $\%2$, $\%4$, $\%5$. Rastgele seçilen bir
ürünün hatalı olma olasılığı?

$B_1, B_2, B_3$: fabrikalara ait ürünler. $A$: hatalı.
\[
  \Prob(A) = 0.02 \times 0.5 + 0.04 \times 0.3 + 0.05 \times 0.2
           = 0.010 + 0.012 + 0.010 = 0.032.
\]
\end{ornek}

\begin{teorem}[Çarpım Kuralı]\label{thm:carpim}
$A_1, \ldots, A_n \in \mathcal{F}$, $\Prob(A_1 \cap \cdots \cap A_{n-1}) > 0$:
\[
  \Prob(A_1 \cap \cdots \cap A_n)
  = \Prob(A_1)\,\Prob(A_2 \mid A_1)\,\Prob(A_3 \mid A_1 \cap A_2)
    \cdots \Prob(A_n \mid A_1 \cap \cdots \cap A_{n-1}).
\]
\end{teorem}

\begin{proof}
Tümevarım: $n=2$ tanımdan. $n$'de geçerliyse $n+1$:
son çarpanı $\Prob(A_1 \cap \cdots \cap A_n) \cdot \Prob(A_{n+1} \mid A_1 \cap \cdots \cap A_n)$
ile genişlet.
\end{proof}

\begin{mikroozet}[§2.3 özeti]
$\Prob(A|B) = \Prob(A \cap B)/\Prob(B)$. Koşullu olasılık bir olasılık ölçüsüdür.
Toplam olasılık: bölüntü üzerinden toplam. Çarpım kuralı: zincirleme hesap.
\end{mikroozet}


% =====================================================================
\section{Bayes Teoremi}
% =====================================================================

\begin{teorem}[Bayes Teoremi]\label{thm:bayes}
$(B_n)$ tam bir bölüntü oluştursun ($\Prob(B_n)>0$, $\bigcup_n B_n = \Omega$,
ikili ayrık). Her $A \in \mathcal{F}$, $\Prob(A)>0$ için:
\[
  \Prob(B_k \mid A)
  = \frac{\Prob(A \mid B_k)\,\Prob(B_k)}{\sum_n \Prob(A \mid B_n)\,\Prob(B_n)}.
\]
\end{teorem}

\begin{proof}
Pay: $\Prob(A \mid B_k)\Prob(B_k) = \Prob(A \cap B_k)$.
Payda: Toplam Olasılık Teoremi (Teorem~\ref{thm:toplam_olasilik}):
$\sum_n \Prob(A \mid B_n)\Prob(B_n) = \Prob(A)$.
Oran: $\Prob(A \cap B_k)/\Prob(A) = \Prob(B_k \mid A)$.
\end{proof}

\begin{notkutu}
Bayes teoremindeki terminoloji:
\begin{itemize}
  \item $\Prob(B_k)$: \textbf{önsel olasılık} \emph{(prior probability)} — $A$ gözlemlenmeden önceki bilgi.
  \item $\Prob(A \mid B_k)$: \textbf{olabilirlik} \emph{(likelihood)} — $B_k$ doğruysa $A$'nın olasılığı.
  \item $\Prob(B_k \mid A)$: \textbf{sonsal olasılık} \emph{(posterior probability)} — $A$ gözlemlendikten sonra.
  \item $\Prob(A)$: \textbf{marjinal (normalleştirme sabiti)}.
\end{itemize}
\end{notkutu}

\begin{ornek}[Bayes --- tıbbi test]\label{ex:tibbi_test}
Hastalığın görülme sıklığı $\Prob(H) = 0.001$. Test duyarlılığı
$\Prob(T^+ \mid H) = 0.99$; yanlış pozitif $\Prob(T^+ \mid H^c) = 0.05$.
Pozitif test sonucunda hasta olma olasılığı?

\[
  \Prob(H \mid T^+)
  = \frac{0.99 \times 0.001}{0.99 \times 0.001 + 0.05 \times 0.999}
  = \frac{0.00099}{0.00099 + 0.04995}
  \approx 0.0194.
\]

Sonuç $\approx \%2$: Hastalık çok nadir olduğunda bile \emph{çok iyi} bir testin
pozitif sonucu gerçek hastalığı büyük olasılıkla değil, yanlış pozitifi gösterir.
\end{ornek}

\begin{hataanalizi}[Savcının yanılgısı (Prosecutor's fallacy)]
$\Prob(E \mid \text{Suçsuz})$ çok küçükse, $\Prob(\text{Suçsuz} \mid E)$'nin
de küçük olduğu \emph{sonucu çıkmaz}. Örnek~\ref{ex:tibbi_test}'te
$\Prob(T^+ \mid H^c) = 0.05$ küçüktür, ama $\Prob(H \mid T^+) \approx 0.02$
daha küçüktür; çünkü $\Prob(H)$ çok küçüktür. Olası suçluluk oranı (önsel)
hesaba katılmadan $\Prob(E \mid \text{Suçsuz})$'u mahkûmiyet gerekçesi yapmak
Bayes teoremini göz ardı etmektir.
\end{hataanalizi}

\begin{disiplinlerarasibaglan}[Makine öğrenmesi ve Naive Bayes]
Bayes teoremi makine öğrenmesinde metin sınıflandırmanın temelini oluşturur.
\emph{Naive Bayes sınıflandırıcısı}, her özelliğin (kelime) bağımsız olduğunu
varsayarak bir belgenin hangi sınıfa ait olduğunu sınıflandırır:
$\Prob(\text{sınıf} \mid \text{kelimeler}) \propto \Prob(\text{sınıf}) \prod_i \Prob(\text{kelime}_i \mid \text{sınıf})$.
``Naive'' (saf) adı bağımsızlık varsayımının kaba olmasından gelir; ama pratikte
şaşırtıcı biçimde iyi çalışır.
\end{disiplinlerarasibaglan}

\begin{mikroozet}[§2.4 özeti]
Bayes teoremi: sonsal $\propto$ olabilirlik $\times$ önsel.
Payda = Toplam Olasılık. Nadir hastalıklar, Savcı yanılgısı, Naive Bayes:
hepsi aynı formülün görünümleri.
\end{mikroozet}


% =====================================================================
\section{Bağımsızlık}
% =====================================================================

\begin{kesif}[Bağımsızlık sezgisi]
``$A$ gerçekleşmesinin $B$'ye etkisi yoktur'' ne demektir?
$\Prob(A \mid B) = \Prob(A)$ formülü bunu yakalar; ama $\Prob(B)=0$ durumunda
bu yazılamaz. Bundan bağımsız bir tanım olasını düşünün.
\ogrencialani[1.5cm]
\end{kesif}

\begin{tanim}[Bağımsız Olaylar]\label{def:bagimsiz_olay}
$A, B \in \mathcal{F}$ olayları \textbf{bağımsız} \emph{(independent)} adını alır;
eğer $\Prob(A \cap B) = \Prob(A)\,\Prob(B)$.

Sonlu aile $\{A_1,\ldots,A_n\}$ \textbf{karşılıklı bağımsız}
\emph{(mutually independent)} adını alır; eğer her $I \subseteq \{1,\ldots,n\}$,
$I \neq \emptyset$ için:
\[
  \Prob\!\left(\bigcap_{i \in I} A_i\right) = \prod_{i \in I} \Prob(A_i).
\]

Sonsuz aile $\{A_i : i \in I\}$ karşılıklı bağımsızdır; eğer her sonlu
$J \subseteq I$ için karşılıklı bağımsızlık koşulu sağlanıyorsa.
\end{tanim}

\begin{kritiksorgu}[İkili bağımsızlık $\neq$ karşılıklı bağımsızlık]
İkili bağımsızlık ($\Prob(A_i \cap A_j) = \Prob(A_i)\Prob(A_j)$ tüm $i \neq j$
için) karşılıklı bağımsızlığı gerektirmez. Aşağıdaki karşı örnek bunu açıklar.
\end{kritiksorgu}

\begin{karsiornek}[Bernstein'ın örneği]\label{ex:bernstein}
Dengeli bir dört yüzlü zar: yüzler $\{1,2,3,4\}$, her biri $1/4$ olasılıklı.
$A = \{1,2\}$, $B = \{1,3\}$, $C = \{1,4\}$ olsun.

$\Prob(A)=\Prob(B)=\Prob(C)=1/2$.
$\Prob(A \cap B) = \Prob(\{1\}) = 1/4 = \Prob(A)\Prob(B)$. Benzer: $A,B$; $A,C$; $B,C$ ikili bağımsız.

Ama $\Prob(A \cap B \cap C) = 1/4 \neq 1/8 = \Prob(A)\Prob(B)\Prob(C)$.

$A, B, C$ ikili bağımsız ama karşılıklı bağımsız değil!
\end{karsiornek}

\begin{tanim}[$\sigma$-Cebirlerinin Bağımsızlığı]\label{def:sigmababag}
$(\mathcal{G}_i)_{i \in I}$ $\mathcal{F}$'nin alt $\sigma$-cebirleri ailesine
\textbf{bağımsız} denir; eğer her sonlu $J \subseteq I$ ve her
$A_j \in \mathcal{G}_j$ ($j \in J$) için:
\[
  \Prob\!\left(\bigcap_{j \in J} A_j\right) = \prod_{j \in J} \Prob(A_j).
\]
\end{tanim}

\begin{onerme}[$\pi$-sistemi ile bağımsızlık kontrolü]\label{prop:pisys_bag}
$\mathcal{P}_i$ bir $\pi$-sistemi ($\cap$ altında kapalı) ve
$\sigma(\mathcal{P}_i) = \mathcal{G}_i$ olsun. Eğer $(\mathcal{P}_i)$ ailesi
bağımsız ise $(\mathcal{G}_i)$ de bağımsızdır.
\end{onerme}

\begin{proof}
Dynkin'in $\pi$-$\lambda$ teoremi: $\pi$-sistemlerin bağımsızlığı,
ürettikleri $\sigma$-cebirlerine uzar.
Ayrıntı için Billingsley~\cite{billingsley1995} §4 bakınız.
\end{proof}

\begin{onerme}[Bağımsızlığın Korunması]\label{prop:bag_korunma}
$A$ ve $B$ bağımsız ise $A$ ve $B^c$; $A^c$ ve $B$; $A^c$ ve $B^c$ de bağımsızdır.
\end{onerme}

\begin{proof}
$\Prob(A \cap B^c) = \Prob(A) - \Prob(A \cap B) = \Prob(A) - \Prob(A)\Prob(B)
= \Prob(A)(1 - \Prob(B)) = \Prob(A)\Prob(B^c)$.
Diğerleri benzer.
\end{proof}

\begin{ornek}[Bağımsız denemeler uzayı]\label{ex:bag_deneme}
Adil yazı-tura $n$ kez atılsın. $\Omega = \{Y,T\}^n$, $|\Omega|=2^n$,
$\Prob(\{\omega\}) = (1/2)^n$ her $\omega$ için.

$A_k = \{k\text{-inci atışta tura}\}$. $A_1, \ldots, A_n$ karşılıklı bağımsızdır:
$\Prob(A_{k_1} \cap \cdots \cap A_{k_r}) = (1/2)^r = \prod_j \Prob(A_{k_j})$.
\end{ornek}

\begin{mikroozet}[§2.5 özeti]
Bağımsızlık: $\Prob(A \cap B) = \Prob(A)\Prob(B)$.
Karşılıklı bağımsızlık: tüm alt aileler için çarpım formülü.
İkili $\not\Rightarrow$ karşılıklı (Bernstein örneği).
Tümleyen bağımsızlığı korur.
\end{mikroozet}


% =====================================================================
\section{Klasik Kombinatoryal Olasılık}
% =====================================================================

\begin{motivasyon}[§2.6]
Sonlu, eşit-olası $\Omega$ ile $\Prob(A) = |A|/|\Omega|$ formülü, olasılık
hesabını sayma problemine indirgeler. Bu alt bölüm önemli sayma tekniklerini
ve klasik paradoksları içeriyor.
\end{motivasyon}

\begin{tanim}[Klasik Olasılık Modeli]\label{def:klasik}
$\Omega$ sonlu, $\mathcal{F} = \mathcal{P}(\Omega)$,
$\Prob(A) = |A|/|\Omega|$ her $A \subseteq \Omega$ için.
\end{tanim}

\subsection{Temel Sayma İlkeleri}

\begin{teorem}[Çarpım İlkesi]\label{thm:carpim_ilkesi}
İlk seçim $m$, ikinci seçim $n$ yolda yapılabiliyorsa, sıralı çift sayısı $mn$'dir.
Genel: $k$ aşamalı seçimde $n_1 \cdots n_k$ olasılık.
\end{teorem}

\begin{tanim}[Permütasyon ve Kombinasyon]\label{def:perm_komb}
\begin{itemize}
  \item $n$ elemanlı kümeden $r$ elemanın \textbf{sıralı} seçimi sayısı:
        $P(n,r) = n(n-1)\cdots(n-r+1) = \dfrac{n!}{(n-r)!}$.
  \item $n$ elemanlı kümeden $r$ elemanın \textbf{sırasız} seçimi sayısı:
        $\binom{n}{r} = \dfrac{n!}{r!\,(n-r)!}$.
  \item $n$ elemanlı kümenin \textbf{tam sıralaması}: $n!$
\end{itemize}
\end{tanim}

\begin{teorem}[Binom Teoremi]\label{thm:binom}
$(x+y)^n = \sum_{k=0}^n \binom{n}{k} x^k y^{n-k}$.
\end{teorem}

\begin{proof}
$(x+y)^n$'de her çarpımı açınca $n$ faktörün her birinden $x$ ya da $y$ seçilir.
$x^k y^{n-k}$ terimini veren seçim sayısı $\binom{n}{k}$ (hangi $k$ faktörden $x$ alınacağı).
\end{proof}

\begin{ornek}[Doğum günü paradoksu]\label{ex:dogumgunu}
$n$ kişilik bir grupta en az iki kişinin aynı doğum gününe sahip olma
olasılığı ne zaman $1/2$'yi aşar?

$A^c$: hiçbir iki kişi aynı günde doğmamış. $|\Omega| = 365^n$,
$|A^c| = 365 \cdot 364 \cdots (365-n+1) = P(365,n)$.
\[
  \Prob(A^c) = \frac{365!/(365-n)!}{365^n} = \prod_{k=0}^{n-1}\frac{365-k}{365}.
\]
$n=23$: $\Prob(A^c) \approx 0.493$; $\Prob(A) \approx 0.507 > 1/2$.

\textbf{Paradoks neden şaşırtıcı?} Biz $n/365$'i düşünürüz; oysa olası çift
sayısı $\binom{n}{2} \approx n^2/2$'dir.
\end{ornek}

\begin{ornek}[Zindan sorunu --- Monty Hall]\label{ex:montyhall}
Üç kapının arkasında: bir araba, iki keçi. Bir kapı seçtiniz; sunucu (keçi
olan) başka bir kapıyı açıyor ve değiştirmenizi öneriyor. Değişmeli misiniz?

$B_k$: arabanın $k$'inci kapıda olması, $k=1,2,3$.
Siz 1'i seçtiniz; sunucu 3'ü açtı.

$\Prob(B_1 \mid \text{3 açıldı}) = \Prob(\text{3 açıldı} \mid B_1)\Prob(B_1) / \Prob(\text{3 açıldı})$.

$\Prob(\text{3 açıldı} \mid B_1) = 1/2$ (sunucu 2 veya 3'ü seçebilir),
$\Prob(\text{3 açıldı} \mid B_2) = 1$ (zorunlu: 3 açılır),
$\Prob(\text{3 açıldı} \mid B_3) = 0$ (3 açılamaz).
$\Prob(\text{3 açıldı}) = 1/2 \cdot 1/3 + 1 \cdot 1/3 = 1/2$.
$\Prob(B_1 \mid \text{3 açıldı}) = (1/6)/(1/2) = 1/3$.
$\Prob(B_2 \mid \text{3 açıldı}) = (1/3)/(1/2) = 2/3$.

\textbf{Değiştirin!} Değiştirme olasılığı $2/3 > 1/3$.
\end{ornek}

\begin{kendinleifade}
Monty Hall sorusunda sunucunun rastgele değil, \emph{her zaman keçi olan}
bir kapıyı açtığını not edin. Sunucu rastgele açsaydı sonuç değişir miydi?
\ogrencialani[2cm]
\end{kendinleifade}

\begin{ornek}[Hat çekme paradoksu]\label{ex:hatchekme}
$n$ kişi arasında $n$ hat dağıtılıyor; $k$ hat kısa ($k < n$). Herkes
rastgele bir hat çekiyor. İlk çeken mi, son çeken mi avantajlı?

$A_j$: $j$-inci kişi kısa hat çekiyor.
$\Prob(A_j) = k/n$ her $j$ için (simetri).

İlk çeken için: $k/n$ açıkça.
İkinci çeken için: $A_1^c$ koşulunda $k/(n-1)$;
$\Prob(A_2) = \Prob(A_2|A_1)\Prob(A_1) + \Prob(A_2|A_1^c)\Prob(A_1^c)$
$= (k-1)/(n-1) \cdot k/n + k/(n-1) \cdot (n-k)/n = k(n-1)/n(n-1) = k/n$.

Hiçbir sıra avantajlı değildir.
\end{ornek}

\begin{mikroozet}[§2.6 özeti]
Klasik model: eşit olası, $\Prob(A) = |A|/|\Omega|$. Çarpım ilkesi, permütasyon,
kombinasyon. Doğum günü: $n=23$'te $>1/2$. Monty Hall: Bayes uygulaması.
Hat çekme: sıranın önemi yok.
\end{mikroozet}


% =====================================================================
\section{Koşullu Bağımsızlık ve Bayesçi Ağlar}
\label{sec:kosullu_bagimsizlik}
% =====================================================================

\begin{tanim}[Koşullu Bağımsızlık]
\label{def:kosullu_bagimsizlik}
$A$ ve $B$ olayları $C$ verildiğinde \textbf{koşullu bağımsız}
(conditionally independent) denir ve $A \perp B \mid C$ gösterilir, eğer
\[
  \Prob(A \cap B \mid C) = \Prob(A \mid C)\,\Prob(B \mid C).
\]
\end{tanim}

\begin{ornek}[Koşullu Bağımsızlık $\neq$ Mutlak Bağımsızlık]
Bir kişi iki gün üst üste hasta olsun ($A$: 1. gün hasta, $B$: 2. gün hasta).
$C$ = "grip mevsiminde". $A$ ve $B$ koşulsuz bağımlı (grip yakalanmak sürer).
$C$ verildiğinde $A$ ve $B$ koşullu bağımsız olabilir: gribi olan kişide
birinci ve ikinci günün durumları hastalığın seyrine bağlı, ama $C$'yi bilirsek
korelasyon büyük ölçüde $C$'den kaynaklanıyor.

Ters yön: $A \perp B$ ama $A \not\perp B \mid C$.
Örnek: $A = \{\text{Türkiye'de yağmur}\}$, $B = \{\text{Japonya'da yağmur}\}$, $C = \{\text{El Niño}\}$.
$A \perp B$ (bağımsız ülkeler) ama $A \not\perp B \mid C$ (El Niño her ikisini etkiler).
\end{ornek}

\begin{teorem}[Koşullu Bağımsızlık Özellikleri]
\label{thm:kosullu_bagimsizlik}
$A \perp B \mid C$ ise:
\begin{enumerate}[(i)]
  \item $A^c \perp B \mid C$ (tümleme kapalı).
  \item $\Prob(A \cap B \cap C) = \Prob(A \mid C)\Prob(B \mid C)\Prob(C)$ (zincir kuralı ile).
  \item $A \perp B \mid C$ ve $A \perp D \mid C$ ise $A \perp (B \cup D) \mid C$ (belirli koşullarla).
\end{enumerate}
\end{teorem}

\begin{proof}[(i)]
$\Prob(A^c \cap B \mid C) = \Prob(B \mid C) - \Prob(A \cap B \mid C)
= \Prob(B \mid C) - \Prob(A \mid C)\Prob(B \mid C) = (1-\Prob(A \mid C))\Prob(B \mid C)
= \Prob(A^c \mid C)\Prob(B \mid C)$.
\end{proof}

\begin{tanim}[Olasılıklı Grafik Model — Bayesçi Ağ]
\label{def:bayesci_ag}
Rasgele değişkenler $X_1, \ldots, X_n$; bir yönlü asiklik graf (DAG).
Her $X_i$, ebeveynleri $\mathrm{Pa}(X_i)$ koşulunda diğer atalardan bağımsız:
\[
  X_i \perp \{X_j : j \notin \mathrm{Pa}(X_i), j \text{ torun değil}\} \mid \mathrm{Pa}(X_i).
\]
Ortak dağılım: $\Prob(X_1 = x_1, \ldots, X_n = x_n) = \prod_{i=1}^n \Prob(X_i = x_i \mid \mathrm{Pa}(X_i))$.
\end{tanim}

\begin{ornek}[Basit Bayesçi Ağ]
$D$ (hastalık) $\to$ $S$ (semptom) $\to$ $T$ (test).
$S \perp T \mid D$ değil, ama $D \perp T \mid S$ (test sadece semptomla ilgili).
Ortak: $\Prob(D, S, T) = \Prob(D)\Prob(S \mid D)\Prob(T \mid S)$.
$\Prob(D \mid T)$: Bayes teoremi ile hesaplanır.
\end{ornek}

\begin{disiplinlerarasibaglan}
\textbf{Yapay zeka:} Bayesçi ağlar, belirsizlik altında akıl yürütmenin
temel matematiksel çerçevesidir; teşhis sistemleri, konuşma tanıma ve robotik karar sistemlerinde kullanılır.\\
\textbf{Epidemiyoloji:} Bulaşıcı hastalık modelleri DAG yapısıyla koşullu bağımsızlıkları kodlar.\\
\textbf{Makine öğrenmesi:} Naive Bayes sınıflandırıcısı, özellikler sınıf koşulunda bağımsız varsayımına dayanır.
\end{disiplinlerarasibaglan}

% =====================================================================
\section{Bilgi Teorisi Temelleri: Entropi ve Karşılıklı Bilgi}
\label{sec:entropi}
% =====================================================================

\begin{tanim}[Shannon entropisi]\label{def:entropi}
Sonlu $\Omega=\{x_1,\ldots,x_n\}$, $p_i=P(X=x_i)$. $X$'in \emph{Shannon entropisi} (\emph{entropy}):
\[
  H(X) = -\sum_{i=1}^n p_i \log_2 p_i,\quad (0\log 0:=0)
\]
bit cinsinden. $H(X)\geq0$; $H(X)=0\iff X$ deterministik; $H(X)\leq\log_2 n$ (maksimum: düzgün dağılım).
\end{tanim}

\begin{teorem}[Maksimum entropi]\label{thm:maks_entropi}
$P(X=x_i)=1/n$ için $H(X)=\log_2 n$; bu $n$-durumlu herhangi bir dağılımın maksimum entropisidir.
\end{teorem}

\begin{proof}
Jensen eşitsizliği: $H(p)=-\sum p_i\log p_i\leq\log_2(\sum p_i/p_i \cdot n^{-1})^{-1}$ ... Doğrudan: Gibbs eşitsizliği, $-\sum p_i\log p_i\leq-\sum p_i\log(1/n)=\log n$, eşitlik $p_i=1/n$.
\end{proof}

\begin{tanim}[Koşullu entropi ve karşılıklı bilgi]\label{def:karsılikli_bilgi}
\begin{align*}
  H(X|Y) &= -\sum_{x,y}P(X=x,Y=y)\log P(X=x|Y=y) \quad \text{(koşullu entropi)}\\
  I(X;Y) &= H(X)-H(X|Y) = H(Y)-H(Y|X) \quad \text{(karşılıklı bilgi)}\\
  D_{\mathrm{KL}}(P\|Q) &= \sum_x P(x)\log\frac{P(x)}{Q(x)} \geq 0 \quad \text{(KL iraksaması)}
\end{align*}
$X\perp Y\iff I(X;Y)=0$.
\end{tanim}

\begin{teorem}[Veri işleme eşitsizliği]\label{thm:veri_isleme}
$X\to Y\to Z$ Markov zinciri ise $I(X;Z)\leq I(X;Y)$: yeterli istatistikler bilgiyi korur, diğer işlemler azaltır.
\end{teorem}

\begin{ornek}[KL iraksaması ile olabilirlik]\label{ex:kl_olasilik}
$P$ gerçek dağılım, $Q_\theta$ model. $D_{\mathrm{KL}}(P\|Q_\theta)=H(P,Q_\theta)-H(P)$; çapraz entropi minimize etmek KL minimize etmekle eşdeğer. MLE = empirik dağılıma en küçük KL olan model.
\end{ornek}

% =====================================================================
\section{Olasılık Uzaylarının Genişlemesi: Sonsuz Dizi Uzayları}
\label{sec:sonsuz_uzay}
% =====================================================================

\begin{tanim}[Ürün uzayı]\label{def:urun_uzayi}
$(\Omega_n,\mathcal{F}_n,P_n)$ olasılık uzayları. Sonsuz ürün $\Omega=\prod_{n=1}^\infty\Omega_n$, $\mathcal{F}=\bigotimes_{n=1}^\infty\mathcal{F}_n$ (silindir kümeler tarafından üretilen $\sigma$-cebir).
\end{tanim}

\begin{teorem}[Kolmogorov genişleme teoremi]\label{thm:kolmogorov_genisleme}
Her tutarlı sonlu boyutlu dağılım ailesi $\{P_{n_1,\ldots,n_k}\}$ için, bu dağılımları marjinal olarak veren tek bir olasılık ölçüsü $P$ sonsuz ürün uzayında mevcuttur.
\end{teorem}

\begin{ornek}[İ.i.d.\ dizi uzayı]\label{ex:iid_uzay}
$X_1,X_2,\ldots\iid F$. Kolmogorov teoremi bu dizinin bir olasılık uzayında var olduğunu garanti eder. Bu, GBSK ve MLT ispatlarının gizli temelidir: sonsuz bir dizi tanımlamak için $(\Omega,\mathcal{F},P)$'nin var olduğunu kabul ederiz.
\end{ornek}

\begin{ispatiskele}
\textbf{Kolmogorov genişleme teoremi ispatı taslağı}
\begin{enumerate}
  \item Sonlu boyutlu dağılımların tutarlılık koşulunu yaz (marjinalizasyon kapalı).
  \item Silindir kümeleri üzerinde $P$'yi tüm tutarlı dağılımları kullanarak tanımla.
  \item Carathéodory genişleme teoremini uygulayarak $\sigma$-cebirde $P$'yi genişlet.
  \item $\sigma$-toplamsallığını (sayılabilir toplanabilirlik) topolojik argümanla doğrula.
\end{enumerate}
\end{ispatiskele}

\begin{mikroozet}
Koşullu bağımsızlık: $\Prob(A\cap B|C) = \Prob(A|C)\Prob(B|C)$.
Koşullu bağımsızlık mutlak bağımsızlığı ne gerektirir ne de garantiler.
Bayesçi ağlar: DAG + koşullu olasılık tabloları $\Rightarrow$ ortak dağılım.
Entropi: belirsizlik ölçüsü; KL iraksaması $\geq0$; MLE = minimum KL.
Kolmogorov genişlemesi: sonsuz i.i.d.\ dizinin varlığı garanti edilir.
\end{mikroozet}

% =====================================================================
%  BÖLÜM SONU ÖZETİ
% =====================================================================
\section*{Bölüm Özeti}
\addcontentsline{toc}{section}{Bölüm Özeti}

\begin{bolumozeti}[Bölüm 2 --- Temel Sonuçlar]
\begin{itemize}
  \item \textbf{Olasılık uzayı:} $(\Omega,\mathcal{F},\Prob)$; (K1) $\Prob \geq 0$;
        (K2) $\Prob(\Omega)=1$; (K3) $\sigma$-toplamsallık.
  \item \textbf{Poincaré:} $\Prob(\bigcup_k A_k) = \sum|I|(-1)^{|I|-1}\Prob(\bigcap_{i\in I}A_i)$.
  \item \textbf{Koşullu olasılık:} $\Prob(A|B) = \Prob(A \cap B)/\Prob(B)$;
        bu bir olasılık ölçüsü.
  \item \textbf{Toplam olasılık:} $\Prob(A) = \sum_n \Prob(A|B_n)\Prob(B_n)$.
  \item \textbf{Bayes:} sonsal $= $ olabilirlik $\times$ önsel / marjinal.
  \item \textbf{Bağımsızlık:} $\Prob(A \cap B) = \Prob(A)\Prob(B)$.
        Karşılıklı bağımsızlık $\neq$ ikili bağımsızlık (Bernstein).
        Tümleyen bağımsızlığı korur.
  \item \textbf{Klasik model:} $\Prob(A) = |A|/|\Omega|$;
        doğum günü $n=23$'te çarpıcı.
\end{itemize}

\medskip\noindent\textbf{Notasyon özeti:}
\begin{itemize}
  \item $\Prob(A|B)$: koşullu olasılık;\quad $\Prob(B_k|A)$: sonsal
  \item $\Prob(B_k)$: önsel;\quad $\Prob(A|B_k)$: olabilirlik
  \item $\binom{n}{k}$: kombinasyon;\quad $n!$: permütasyon
\end{itemize}
\end{bolumozeti}

\begin{ozyansima}
\begin{itemize}
  \item Kolmogorov aksiyomlarından tüm temel özellikleri türetebilir misiniz
        (tümleme, monotonluk, Poincaré)?
  \item Koşullu olasılığın neden bir olasılık ölçüsü olduğunu üç aksiyomu
        kontrol ederek açıklayabilir misiniz?
  \item Tıbbi test örneğini farklı bir hastalık sıklığı ($1/100$) ile yeniden
        hesaplayın; sonuç nasıl değişti?
  \item Bernstein'ın karşı örneğinde neden karşılıklı bağımsızlık bozuluyor?
\end{itemize}
\end{ozyansima}

\begin{kopru}
\textbf{Sonraki bölüm (Bölüm~3: Rasgele Değişkenler ve Dağılımlar).}
Olasılık uzayı altyapısı artık hazır. Rasgele değişken, $(\Omega,\mathcal{F})$'den
$(\mathbb{R},\mathcal{B}(\mathbb{R}))$'ye ölçülebilir fonksiyon olarak
tanımlanacak; dağılım fonksiyonu bu fonksiyonun $\Prob$ altındaki görüntüsü
olacak.

\medskip

\textbf{İleri okuma.}
Kolmogorov~\cite{kolmogorov1933} (orijinal kaynak, 77 sayfa);
Feller~\cite{feller1968vol1} Bölüm~I--V;
Gnedenko~\cite{gnedenko1962} Bölüm~1--3.
\end{kopru}


% =====================================================================
%  ALIŞTIRMALAR
% =====================================================================

\cozumlualistirmalar

\begin{alistirma}[Al.2.1]{A}{Poincaré formülü --- üç olay}
$\Prob(A) = 0.5$, $\Prob(B) = 0.4$, $\Prob(C) = 0.3$,
$\Prob(A \cap B) = 0.2$, $\Prob(A \cap C) = 0.15$,
$\Prob(B \cap C) = 0.1$, $\Prob(A \cap B \cap C) = 0.05$.
$\Prob(A \cup B \cup C)$'yi hesaplayın.
\end{alistirma}
\alcozum{%
Poincaré (Teorem~\ref{thm:poincare}):
\begin{align*}
  \Prob(A \cup B \cup C)
  &= 0.5 + 0.4 + 0.3 - 0.2 - 0.15 - 0.1 + 0.05 = 0.80.
\end{align*}
$\blacksquare$%
}

\begin{alistirma}[Al.2.2]{B}{Bayes teoremi --- fabrika}
Örnek~\ref{ex:fabrika}'ya devam: Seçilen ürün hatalı çıktı.
Hangi fabrikadan geldiğinin sonsal olasılıklarını hesaplayın.
\end{alistirma}
\alcozum{%
$\Prob(A) = 0.032$ (Örnek~\ref{ex:fabrika}'dan).
\[
  \Prob(B_1 \mid A) = \frac{0.02 \times 0.5}{0.032} = \frac{0.010}{0.032} \approx 0.3125.
\]
\[
  \Prob(B_2 \mid A) = \frac{0.04 \times 0.3}{0.032} = \frac{0.012}{0.032} = 0.375.
\]
\[
  \Prob(B_3 \mid A) = \frac{0.05 \times 0.2}{0.032} = \frac{0.010}{0.032} \approx 0.3125.
\]
Toplam: $0.3125 + 0.375 + 0.3125 = 1$. \checkmark
Hatalı ürün en olası fabrika 2'den geliyor. $\blacksquare$%
}

\begin{alistirma}[Al.2.3]{A}{Bağımsızlık doğrulama}
Adil zar: $\Omega = \{1,2,3,4,5,6\}$, $\Prob(\{k\}) = 1/6$.
$A = \{1,2,3,4\}$, $B = \{3,4,5,6\}$, $C = \{1,3,5\}$.
(a) $A$ ve $B$ bağımsız mı?
(b) $A$ ve $C$ bağımsız mı?
(c) $B$ ve $C$ bağımsız mı?
(d) $A, B, C$ karşılıklı bağımsız mı?
\end{alistirma}
\alcozum{%
$\Prob(A) = 4/6 = 2/3$, $\Prob(B) = 4/6 = 2/3$, $\Prob(C) = 3/6 = 1/2$.

(a) $A \cap B = \{3,4\}$; $\Prob(A \cap B) = 2/6 = 1/3$.
$\Prob(A)\Prob(B) = 4/9 \neq 1/3$. Bağımsız \emph{değil}.

(b) $A \cap C = \{1,3\}$; $\Prob(A \cap C) = 2/6 = 1/3$.
$\Prob(A)\Prob(C) = 2/3 \cdot 1/2 = 1/3$. Bağımsız. \checkmark

(c) $B \cap C = \{3,5\}$; $\Prob(B \cap C) = 1/3$.
$\Prob(B)\Prob(C) = 1/3$. Bağımsız. \checkmark

(d) $A \cap B \cap C = \{3\}$; $\Prob = 1/6$.
$\Prob(A)\Prob(B)\Prob(C) = 2/3 \cdot 2/3 \cdot 1/2 = 2/9 \neq 1/6$.
Karşılıklı bağımsız \emph{değil}. $\blacksquare$%
}

\begin{alistirma}[Al.2.4]{B}{Doğum günü --- genelleştirilmiş}
$n = 50$ kişilik bir grupta aynı doğum gününe sahip iki kişinin olmama
olasılığını hesaplayın (logaritma kullanın). $\Prob(A)$ nedir?
\end{alistirma}
\alcozum{%
$\Prob(A^c) = \prod_{k=0}^{49}\frac{365-k}{365}$.
$\ln \Prob(A^c) = \sum_{k=0}^{49} \ln(1 - k/365)$.
$\ln(1-k/365) \approx -k/365$ için kaba tahmin:
$\sum_{k=0}^{49} k/365 = \binom{50}{2}/365 = 1225/365 \approx 3.356$.
$\Prob(A^c) \approx e^{-3.356} \approx 0.035$.
$\Prob(A) \approx 0.965$.

Daha kesin: $\Prob(A^c) \approx 0.0296$; $\Prob(A) \approx 0.970$. $\blacksquare$%
}

\begin{alistirma}[Al.2.4b]{B}{Koşullu Olasılık Paradoksu: Monty Hall}
Monty Hall problemi: 3 kapı, birinde araba diğer ikisinde keçi. Yarışmacı kapı 1'i seçer. Sunucu, keçili bir kapıyı açar (kapı 3). Değiştirmeli mi?
(a) $A$: kapı 1'de araba; $B$: sunucu kapı 3'ü açar; Bayes teoremi ile $P(A|B)$'yi hesaplayın.
(b) $P(\text{araba}|\text{değiştirir})$'yi bulun.
(c) Sezgisiz insanların neden yanıldığını olasılık teorisi açısından açıklayın.
\end{alistirma}
\alcozum{%
(a) Toplam olasılık: $P(B) = P(B|A)P(A)+P(B|A^c)P(A^c)$.
$P(A)=1/3$; $P(A^c)=2/3$.
$P(B|A)=1/2$ (sunucu 2 veya 3'ü eşit olasılıkla açar; biz kapı 3'ü koşullandırdık).
$P(B|A^c)=1$ eğer araba kapı 2'deyse sunucu sadece kapı 3'ü açabilir; $P(B|A^c=\text{kapı 2})=1$, $P(B|A^c=\text{kapı 3})=0$.

Daha dikkatli: $A_k$: araba kapı $k$'da.
$P(A_1)=P(A_2)=P(A_3)=1/3$.
Sunucu kapı 3'ü açar: $P(B|A_1)=1/2$, $P(B|A_2)=1$, $P(B|A_3)=0$.
$P(B)=\frac{1}{3}\cdot\frac{1}{2}+\frac{1}{3}\cdot1+\frac{1}{3}\cdot0=\frac{1}{2}$.
$P(A_1|B)=\frac{(1/2)(1/3)}{1/2}=\frac{1}{3}$; $P(A_2|B)=\frac{1\cdot(1/3)}{1/2}=\frac{2}{3}$.

(b) Değiştirir $\to$ kapı 2 seçer: $P(\text{araba})=P(A_2|B)=2/3$. Değiştirmez: $1/3$. Değiştirmek iki kat daha iyi!

(c) Yanılgı: Sunucunun ek bilgi verdiği göz ardı ediliyor. Sunucunun kapı açması rastgele değil — keçi içeren kapıyı bilinçli açar. Bu bilgi kapı 2'nin olasılığını günceller. İnsanlar genellikle sunucunun hamlesi bağımsız diye düşünür ($P=1/3$ değişmez sanır); Bayes teoremi bu bağımlılığı doğru biçimde yakalar.}

\begin{alistirma}[Al.2.4c]{C}{Olasılık Ölçüsünün Süreklilik Özellikleri}
$(\Omega,\mathcal{F},\Prob)$ olasılık uzayı.
(a) $A_n\uparrow A$ (artan limit) için $\Prob(A_n)\to\Prob(A)$ olduğunu kanıtlayın.
(b) $B_n\downarrow B$ (azalan limit) için $\Prob(B_n)\to\Prob(B)$ olduğunu kanıtlayın.
(c) $\Prob(\limsup_n A_n)$'ı yorumlayın ve Borel-Cantelli ile bağlantısını kurun.
\end{alistirma}
\alcozum{%
(a) \textbf{Arttan süreklilik}: $C_1=A_1$, $C_n=A_n\setminus A_{n-1}$ (ayrık) olsun. $A=\bigcup_n C_n$ (ayrık birleşim).
$\Prob(A)=\sum_{k=1}^\infty\Prob(C_k)=\lim_{n\to\infty}\sum_{k=1}^n\Prob(C_k)=\lim_{n\to\infty}\Prob(A_n)$.
(Sayılabilir toplamsellik ve dizi yakınsaması birlikte.)

(b) \textbf{Azalan süreklilik}: $D_n=B_1\setminus B_n\uparrow B_1\setminus B$.
(a)'dan $\Prob(D_n)\to\Prob(B_1\setminus B)$.
$\Prob(D_n)=\Prob(B_1)-\Prob(B_n)$ (tek yönlü kapsama); $\Prob(B_n)\to\Prob(B)$. $\checkmark$

(c) $\limsup_n A_n=\{A_n\text{ sonsuz sık gerçekleşir}\}=\bigcap_{n=1}^\infty\bigcup_{k=n}^\infty A_k$.
Borel-Cantelli 1: $\sum\Prob(A_n)<\infty\Rightarrow\Prob(\limsup A_n)=0$ (sonsuz sık gerçekleşmez).
Borel-Cantelli 2: $A_n$ bağımsız, $\sum\Prob(A_n)=\infty\Rightarrow\Prob(\limsup A_n)=1$ (sonsuz sık gerçekleşir).
Bu iki yönlü sonuç, neredeyse kesin yakınsama kriterleri için temel araçtır.}

\alistirmalar

\begin{alistirma}[Al.2.5]{A}{Olasılık aksiyomlarından türetme}
$(\Omega,\mathcal{F},\Prob)$ olasılık uzayı, $A, B \in \mathcal{F}$.
Yalnızca Kolmogorov aksiyomlarını (K1)--(K3) kullanarak gösterin:
\begin{enumerate}[label=(\alph*)]
  \item $\Prob(A) \leq 1$.
  \item $\Prob(A^c) = 1 - \Prob(A)$.
  \item $A \subseteq B \Rightarrow \Prob(A) \leq \Prob(B)$.
\end{enumerate}
\end{alistirma}
\alcozum{%
(a) $\Omega = A \cup A^c$ ayrık; (K3) + (K1): $1 = \Prob(A)+\Prob(A^c) \geq \Prob(A)$.
(b) Aynı ayrışım: $\Prob(A^c) = 1 - \Prob(A)$.
(c) $B = A \cup (B \setminus A)$ ayrık; $\Prob(B) = \Prob(A)+\Prob(B \setminus A) \geq \Prob(A)$. $\blacksquare$%
}

\begin{alistirma}[Al.2.6]{A}{Koşullu olasılık zinciri}
Bir kutuda 3 kırmızı, 4 mavi top var. Geri koymadan iki top çekiliyor.
(a) İlk kırmızı, ikinci mavi çekme olasılığı.
(b) Her ikisi de mavi çekme olasılığı.
(c) En az bir kırmızı çekme olasılığı.
\end{alistirma}
\alcozum{%
(a) $\Prob(K_1 \cap M_2) = \Prob(K_1)\Prob(M_2|K_1) = \frac{3}{7} \cdot \frac{4}{6} = \frac{12}{42} = \frac{2}{7}$.

(b) $\Prob(M_1 \cap M_2) = \frac{4}{7} \cdot \frac{3}{6} = \frac{12}{42} = \frac{2}{7}$.

(c) $\Prob(\text{en az bir kırmızı}) = 1 - \Prob(M_1 \cap M_2) = 1 - 2/7 = 5/7$. $\blacksquare$%
}

\begin{alistirma}[Al.2.7]{A}{Boole eşitsizliği}
$n$ olay için Boole eşitsizliğini doğrulayın:
$\Prob\!\bigl(\bigcap_{k=1}^n A_k\bigr) \geq 1 - \sum_{k=1}^n \Prob(A_k^c)$.
Bu eşitsizliğin adı Bonferroni eşitsizliğidir.
\end{alistirma}
\alcozum{%
$\Prob\!\bigl(\bigcap_k A_k\bigr) = 1 - \Prob\!\bigl(\bigcup_k A_k^c\bigr)$
(De~Morgan ve tümleyen kuralı).
Alt-toplamsallık: $\Prob(\bigcup_k A_k^c) \leq \sum_k \Prob(A_k^c)$.
Dolayısıyla $\Prob(\bigcap_k A_k) \geq 1 - \sum_k \Prob(A_k^c)$. $\blacksquare$%
}

\begin{alistirma}[Al.2.8]{B}{Koşullu bağımsızlık}
$A$ ve $B$, $C$'ye koşullu bağımsız adını alır; eğer
$\Prob(A \cap B \mid C) = \Prob(A \mid C)\Prob(B \mid C)$.
(a) $A$ ve $B$ koşulsuz bağımsızsa, $C$'ye koşullu bağımsız olması gerekir mi?
    Karşı örnek verin.
(b) $A$ ve $B$, $C$'ye koşullu bağımsızsa, koşulsuz bağımsız olması gerekir mi?
    Karşı örnek verin.
\end{alistirma}
\alcozum{%
(a) Hayır. İki adil zar: $A =$ \{1. zar çift\}, $B =$ \{2. zar çift\}, $C = \{$toplam 7$\}$.
$A,B$ bağımsız. Ama $\Prob(A|C) = 3/6 = 1/2$ ve $\Prob(B|C) = 3/6 = 1/2$;
$\Prob(A \cap B \mid C) = \Prob(\{(2,5),(4,3),(6,1)\} \mid C) = 0$
(çift+çift=çift, toplam 7 olamaz). $\neq 1/4$. Koşullu bağımsız değil.

(b) Hayır. Örnek: $\Omega=\{1,2,3,4\}$ düzgün; $A=\{1,2\}$, $B=\{2,3\}$, $C=\{1,3\}$.
$\Prob(A)=\Prob(B)=1/2$, $\Prob(A\cap B)=1/4=\Prob(A)\Prob(B)$: bağımsız.
$\Prob(A\cap B\mid C)=\Prob(\emptyset\mid C)=0$,
$\Prob(A|C)\Prob(B|C)=\frac{1}{2}\cdot\frac{1}{2}=\frac{1}{4}\neq 0$:
koşullu bağımsız değil. $\blacksquare$%
}

\begin{alistirma}[Al.2.9]{B}{Toplam olasılık --- iki aşamalı deney}
Bir torba 5 beyaz, 3 siyah top içeriyor. Geri koymadan iki top çekiliyor.
$A_k =$ \{$k$-inci top beyaz\}.
(a) $\Prob(A_1)$, $\Prob(A_2)$, $\Prob(A_1 \cap A_2)$'yi hesaplayın.
(b) $\Prob(A_2)$'yi Toplam Olasılık Teoremi ile doğrulayın.
(c) $\Prob(A_1 \mid A_2)$ ve $\Prob(A_2 \mid A_1)$'i hesaplayın.
\end{alistirma}
\alcozum{%
(a) $\Prob(A_1) = 5/8$.
$\Prob(A_2) = \Prob(A_2|A_1)\Prob(A_1)+\Prob(A_2|A_1^c)\Prob(A_1^c)$
$= \frac{4}{7}\cdot\frac{5}{8}+\frac{5}{7}\cdot\frac{3}{8} = \frac{20}{56}+\frac{15}{56}=\frac{35}{56}=\frac{5}{8}$.
$\Prob(A_1 \cap A_2) = \frac{5}{8}\cdot\frac{4}{7} = \frac{20}{56} = \frac{5}{14}$.

(b) Yukarıdaki hesap Toplam Olasılık Teoremi'nin uygulaması.

(c) $\Prob(A_1|A_2) = \Prob(A_1\cap A_2)/\Prob(A_2) = (5/14)/(5/8) = 8/14 = 4/7$.
$\Prob(A_2|A_1) = \Prob(A_1\cap A_2)/\Prob(A_1) = (5/14)/(5/8) = 4/7$. $\blacksquare$%
}

\begin{alistirma}[Al.2.10]{B}{Hat çekme genelleştirilmiş}
$n$ kişi, $m$ şanslı hat ($m < n$). Hiç kimse önceki seçimi görmeden
kendi hatını çekiyor (tüm sıralamalar eşit olası).
$j$-inci kişinin şanslı hat çekme olasılığının $j$'den bağımsız olduğunu
tümevarımla ispatlayın.
\end{alistirma}
\alcozum{%
$\Prob(A_j) = m/n$ her $j$ için.

$j=1$: $\Prob(A_1) = m/n$ açıkça.

$j$'de doğru varsayalım; $j+1$:
$\Prob(A_{j+1}) = \Prob(A_{j+1}|A_j)\Prob(A_j) + \Prob(A_{j+1}|A_j^c)\Prob(A_j^c)$
$= \frac{m-1}{n-1}\cdot\frac{m}{n} + \frac{m}{n-1}\cdot\frac{n-m}{n}$
$= \frac{m(m-1) + m(n-m)}{n(n-1)} = \frac{m(n-1)}{n(n-1)} = \frac{m}{n}$.
$\blacksquare$%
}

\begin{alistirma}[Al.2.11]{C}{$\pi$-$\lambda$ teoremi ve bağımsızlık}
$\mathcal{P}_1 = \{A_1 \times A_2 : A_i \in \mathcal{F}_i\}$ ikili ayrık
dikdörtgenler ailesi; $\mathcal{P}_1$ bir $\pi$-sistemdir.
$\Prob_1$ ve $\Prob_2$ aynı $\pi$-sistem üzerinde aynı değerleri alıyorsa,
$\sigma(\mathcal{P}_1)$ üzerinde eşit olduklarını Dynkin'in $\pi$-$\lambda$
teoremini kullanarak ispatlayın. Bu teoremin ürün ölçüsünün tekliği ile
bağlantısını açıklayın.
\end{alistirma}
\alcozum{%
\begin{ipucu}{1}
Dynkin'in $\pi$-$\lambda$ teoremi: $\mathcal{P}$ bir $\pi$-sistem (sonlu kesişim
altında kapalı) ve $\lambda(\mathcal{P}) = \mathcal{P}$ içeren en küçük
Dynkin sistemi ise $\lambda(\mathcal{P}) = \sigma(\mathcal{P})$.
\end{ipucu}

$D = \{B \in \sigma(\mathcal{P}_1) : \Prob_1(B) = \Prob_2(B)\}$ tanımlayın.
$D$ bir Dynkin sistemidir: (i) $\Omega \in D$; (ii) $A \subseteq B$, $A,B \in D
\Rightarrow B \setminus A \in D$; (iii) $A_n \nearrow A$, $A_n \in D \Rightarrow A \in D$.

$\mathcal{P}_1 \subseteq D$ (verilen koşul).
Dynkin teoreminden $\sigma(\mathcal{P}_1) \subseteq D$, yani $\Prob_1 = \Prob_2$ üzerinde.

Ürün ölçüsü: $\mu_1 \otimes \mu_2$ dikdörtgenlerde
$\mu_1(A_1)\mu_2(A_2)$'yi üretir; bu değerler başka bir ölçüyü belirliyorsa
her ikisi $\sigma(\mathcal{P}_1)$ üzerinde eşittir; yani ürün ölçüsü
$\sigma$-sonlu durumlarda tekdir. $\blacksquare$%
}

\begin{alistirma}[Al.2.12]{C}{Sonsuz bağımsız denemeler}
$(\Omega, \mathcal{F}, \Prob)$ üzerinde $A_1, A_2, \ldots$ bağımsız,
$\Prob(A_n) = p \in (0,1)$ her $n$ için.
(a) $\Prob(A_n \text{ s.s. çok})$ — yani $\Prob(\limsup_n A_n)$ — nedir?
(b) $\Prob(A_n \text{ eninde sonunda hep gerçekleşir})$ nedir?
(c) Cevaplarınızı $p$ değerine göre yorumlayın.
\end{alistirma}
\alcozum{%
(a) $\sum_n \Prob(A_n) = \sum_n p = \infty$ (sonsuz toplam).
Borel-Cantelli II (Bölüm~7'de kanıtlanacak): $A_n$ bağımsız ve $\sum \Prob(A_n) = \infty$
ise $\Prob(\limsup_n A_n) = 1$.

(b) $\liminf_n A_n$: eninde sonunda hep gerçekleşir.
$B_n = \bigcap_{k \geq n} A_k^c$; $\Prob(B_n) = \prod_{k \geq n}(1-p) = 0$
(sonsuz çarpım, $p>0$). $\Prob(\liminf_n A_n) = 1 - \Prob(\limsup_n A_n^c)$.
$\sum_n \Prob(A_n^c) = \infty$; $A_n^c$ da bağımsız; Borel-Cantelli II:
$\Prob(\limsup_n A_n^c) = 1$.
$\Prob(\liminf_n A_n) = 0$.

(c) Her $p \in (0,1)$ için: sonsuz bağımsız denemede her olay sonsuz sık
gerçekleşir (a.s.), ama hiçbir zaman sonunda \emph{sürekli} gerçekleşmez. $\blacksquare$%
}

\begin{alistirma}[Al.2.13]{A}{Koşullu Bağımsızlık}
Sağlık veri tabanında: $D$ = hastada kanser var, $S$ = semptom pozitif, $T$ = test pozitif.
Verilen: $\Prob(D) = 0.01$, $\Prob(S \mid D) = 0.8$, $\Prob(S \mid D^c) = 0.05$,
$\Prob(T \mid S) = 0.9$, $\Prob(T \mid S^c) = 0.1$.
Bayesçi ağ: $D \to S \to T$.
(a) $\Prob(S)$'yi hesaplayın. (b) $\Prob(T)$'yi hesaplayın. (c) $\Prob(D \mid T)$'yi hesaplayın.
\end{alistirma}
\alcozum{%
(a) $\Prob(S) = 0.8\cdot0.01 + 0.05\cdot0.99 = 0.008 + 0.0495 = 0.0575$.\\
(b) $\Prob(T) = 0.9\cdot\Prob(S) + 0.1\cdot\Prob(S^c) = 0.9(0.0575) + 0.1(0.9425) = 0.05175 + 0.09425 = 0.146$.\\
(c) $\Prob(D \mid T) = \Prob(T \mid D)\Prob(D)/\Prob(T)$.
$\Prob(T \mid D) = 0.9\cdot0.8 + 0.1\cdot0.2 = 0.74$.
$\Prob(D \mid T) = 0.74\cdot0.01/0.146 \approx 0.051$. Yüksek test duyarlılığına rağmen düşük: nadir hastalık etkisi.}

\begin{alistirma}[Al.2.14]{B}{Koşullu Bağımsızlık vs Mutlak Bağımsızlık}
$\Omega = \{1,2,3,4\}$, eşit olası.
$A = \{1,2\}$, $B = \{1,3\}$, $C = \{1,4\}$.
(a) $A$ ve $B$ bağımsız mı? (b) $A$ ve $B$, $C$ koşulunda bağımsız mı?
(c) $A$ ve $C$ bağımsız mı? $A$ ve $C$, $B$ koşulunda?
\end{alistirma}
\alcozum{%
$\Prob(A) = \Prob(B) = \Prob(C) = 1/2$.
$\Prob(A\cap B) = \Prob(\{1\}) = 1/4 = \Prob(A)\Prob(B)$: (a) evet, $A \perp B$.\\
(b) $\Prob(C) = 1/2$; $\Prob(A\cap B\cap C) = \Prob(\{1\}) = 1/4$.
$\Prob(A\cap B \mid C) = \Prob(A\cap B\cap C)/\Prob(C) = (1/4)/(1/2) = 1/2$.
$\Prob(A\mid C)\Prob(B\mid C) = (1/2)(1/2) = 1/4 \neq 1/2$.
Koşullu bağımsız değil!\\
(c) $A \perp C$: $\Prob(A\cap C) = 1/4 = \Prob(A)\Prob(C)$. Evet.
$\Prob(A\cap C\mid B) = \Prob(A\cap B\cap C)/\Prob(B) = (1/4)/(1/2) = 1/2 \neq \Prob(A|B)\Prob(C|B) = (1/2)(1/2)$.
$A \not\perp C \mid B$.}

\begin{alistirma}[Al.2.15]{B}{Bayesçi Ağ — Ortak Dağılım}
$X_1$, $X_2 \mid X_1$, $X_3 \mid X_2$ zinciri (her ikili normal):
$X_1 \sim \Normal(0,1)$; $X_2 \mid X_1 = x_1 \sim \Normal(x_1, 1)$; $X_3 \mid X_2 = x_2 \sim \Normal(x_2, 1)$.
(a) $E[X_2]$ ve $\mathrm{Var}(X_2)$'yi hesaplayın.
(b) $X_1$ ve $X_3$'ün koşulsuz korelasyonunu bulun.
(c) $X_3 \mid X_2 = 0$ ile $X_3 \mid X_1 = 0$'ı karşılaştırın.
\end{alistirma}
\alcozum{%
(a) $E[X_2] = E[E[X_2|X_1]] = E[X_1] = 0$.
$\mathrm{Var}(X_2) = E[\mathrm{Var}(X_2|X_1)] + \mathrm{Var}(E[X_2|X_1]) = 1 + \mathrm{Var}(X_1) = 1+1 = 2$.\\
(b) Aynı şekilde: $E[X_3] = 0$, $\mathrm{Var}(X_3) = 3$.
$\mathrm{Cov}(X_1,X_3) = E[X_1 X_3]$. Kule beklentisi: $E[X_1 X_3] = E[X_1 E[X_3|X_1,X_2]]$... zincirde: $E[X_3|X_2] = X_2$; $E[X_1 X_3] = E[X_1 X_2] = E[X_1^2] = 1$.
$\mathrm{Cor}(X_1,X_3) = 1/\sqrt{1\cdot3} = 1/\sqrt{3} \approx 0.577$.\\
(c) $X_3|X_2=0 \sim \Normal(0,1)$. $X_3|X_1=0$: $X_2|X_1=0 \sim \Normal(0,1)$, sonra $X_3|X_2 \sim \Normal(X_2,1)$; marjinal $X_3|X_1=0 \sim \Normal(0,2)$.
Daha geniş: $X_1$ bilgisi $X_2$ üzerinden dolaylı.}
